% Informe técnico — TPI L-Cero (formato de la plantilla UBP)
% Compilar desde docs/informe/:   latexmk -xelatex informe.tex
% Las tablas de generado/ se regeneran con: python scripts/generar_documentacion.py

% ------------------------------------------------------------------ COMPLETAR
\newcommand{\Asignatura}{[ASIGNATURA]}
\newcommand{\Docente}{RODRÍGUEZ ISLEÑO, SANTIAGO HORACIO}
\newcommand{\EtiquetaDocente}{Docente}            % p. ej. Profesora, Profesor
\newcommand{\Autor}{RODRÍGUEZ, RODRIGO EMMANUEL \\ CARRIZO, MATHIAS EZEQUIEL \\ RODRÍGUEZ, TOMÁS AGUSTÍN}  % uno por línea, separados con \\
\newcommand{\EtiquetaAutor}{Autores}               % Autor, si es uno solo
\newcommand{\Repositorio}{[URL del repositorio Git de la cátedra]}
\newcommand{\Titulo}{“L-Cero: compilador y máquina virtual abstracta”}
\newcommand{\Actividad}{Evaluación Parcial N°2 (Trabajo Práctico Integrador)}
\newcommand{\FechaCorta}{02/11/2026}
\newcommand{\FechaLarga}{Córdoba, 2 de noviembre de 2026}
\newcommand{\Carrera}{Ing. en Informática}
% -----------------------------------------------------------------------------

\documentclass[11pt,a4paper]{article}

% --- Página: mismas medidas que la plantilla de Word (márgenes de 3,17 y 2,54 cm)
\usepackage[a4paper, left=90pt, right=72.9pt, top=131pt, bottom=82pt,
            headheight=12pt, headsep=8pt, footskip=20pt]{geometry}

% --- Tipografías: Liberation Sans tiene las mismas medidas que Arial y Carlito que Calibri
\usepackage{fontspec}
\usepackage{polyglossia}
\setdefaultlanguage{spanish}
\setmainfont{Liberation Sans}
\newfontfamily\calibri{Carlito}
\setmonofont{DejaVu Sans Mono}[Scale=0.8]
\linespread{0.97}                         % interlineado de 13,2 pt, como en Word
\setlength{\parindent}{25.6pt}
\setlength{\parskip}{7pt plus 1pt}
\hyphenpenalty=10000 \exhyphenpenalty=10000   % Word no corta palabras
\setlength{\emergencystretch}{4em}

\usepackage{amsmath,amssymb,amsthm}
\usepackage{graphicx,float}
\usepackage{array,longtable,makecell}
\usepackage[shortlabels]{enumitem}
\usepackage{fvextra}
\usepackage{xcolor}
\usepackage{tikz}
\usetikzlibrary{arrows.meta,positioning}
\usepackage{pifont}
\usepackage{lastpage}
\usepackage{fancyhdr}
\usepackage{titlesec}
\usepackage{tocloft}
\usepackage[hidelinks]{hyperref}
\usepackage[font=it, labelfont=it, labelsep=period, justification=centering, skip=5pt]{caption}

\gappto\captionsspanish{%
  \renewcommand{\tablename}{Tabla}%
  \renewcommand{\figurename}{Figura}%
  \renewcommand{\contentsname}{ÍNDICE}}

% --- Títulos (Calibri negrita, como en la plantilla)
\setcounter{secnumdepth}{2}
\titleformat{\section}{\calibri\bfseries\fontsize{13}{15.6}\selectfont}{\thesection.}{0.28em}{}
\titleformat{\subsection}{\calibri\bfseries\fontsize{10.5}{12.6}\selectfont}{\thesubsection}{0.28em}{}
\titleformat{\subsubsection}{\calibri\bfseries\fontsize{10}{12}\selectfont}{}{0pt}{}
\titlespacing*{\section}{0pt}{10pt}{2pt}
\titlespacing*{\subsection}{0pt}{6pt}{0pt}
\titlespacing*{\subsubsection}{0pt}{4pt}{0pt}

% --- Índice plano con puntos, como el de Word
\renewcommand{\cfttoctitlefont}{\calibri\bfseries\fontsize{12}{14.4}\selectfont}
\renewcommand{\cftaftertoctitle}{}
\setlength{\cftbeforetoctitleskip}{0pt}
\setlength{\cftaftertoctitleskip}{6pt}
\renewcommand{\cftsecfont}{\normalfont}
\renewcommand{\cftsecpagefont}{\normalfont}
\renewcommand{\cftsecleader}{\cftdotfill{0.6}}
\renewcommand{\cftsubsecleader}{\cftdotfill{0.6}}
\renewcommand{\cftsecaftersnum}{.}
\setlength{\cftsecindent}{0pt}
\setlength{\cftsubsecindent}{0pt}
\setlength{\cftbeforesecskip}{2pt}
\setlength{\cftbeforesubsecskip}{2pt}
\tocloftpagestyle{ubp}
\settowidth{\cftsecnumwidth}{9.\ }
\settowidth{\cftsubsecnumwidth}{9.9\ }
\cftsetpnumwidth{1em}
\cftsetrmarg{1.5em}

% --- Encabezado y pie de la plantilla UBP (posiciones medidas sobre el original)
\newcommand{\encabezadoUBP}{%
\begin{tikzpicture}[remember picture, overlay, x=1pt, y=-1pt, shift={(current page.north west)}]
  \node[anchor=north west, inner sep=0] at (90,41.8) {\includegraphics[width=58pt]{plantilla/ubp_logo_encabezado.png}};
  \node[anchor=base, inner sep=0, font=\bfseries\fontsize{11}{13}\selectfont] at (311,68.7) {UNIVERSIDAD BLAS PASCAL};
  \node[anchor=base, inner sep=0, font=\fontsize{11}{13}\selectfont] at (311,81.4) {MODALIDAD ON CAMPUS};
  \node[anchor=base east, inner sep=0, font=\fontsize{11}{13}\selectfont] at (522.4,75.1)
    {{\fontsize{8}{9}\selectfont Página} \thepage{} {\fontsize{8}{9}\selectfont de} \pageref*{LastPage}};
  \draw[line width=0.75pt] (90,113) -- (522.4,113);
  \draw[line width=0.9pt] (90,787) -- (522.4,787);
  \node[anchor=base west, inner sep=0, font=\fontsize{9}{11}\selectfont] at (90,797.8) {Fecha: \FechaCorta};
  \node[anchor=base east, inner sep=0, font=\fontsize{9}{11}\selectfont] at (522.4,797.8) {\Carrera};
\end{tikzpicture}}
\fancypagestyle{ubp}{%
  \fancyhf{}%
  \renewcommand{\headrulewidth}{0pt}%
  \renewcommand{\footrulewidth}{0pt}%
  \fancyhead[C]{\encabezadoUBP}}
\pagestyle{ubp}

% --- Tablas con grilla, como las de Word
\renewcommand{\theadfont}{\bfseries}
\newcolumntype{L}[1]{>{\raggedright\arraybackslash}m{#1}}
\newcolumntype{C}[1]{>{\centering\arraybackslash}m{#1}}
\newcommand{\toprule}{\hline}
\newcommand{\midrule}{\hline}
\newcommand{\bottomrule}{\hline}
\newenvironment{grilla}{\footnotesize\renewcommand{\arraystretch}{1.45}\setlength{\tabcolsep}{5pt}}{}

% --- Teoremas (sin cursiva, como se escribirían en Word)
\theoremstyle{definition}
\newtheorem{teorema}{Teorema}
\newtheorem{definicion}{Definición}

% --- Notación
\newcommand{\nt}[1]{\textit{#1}}
\newcommand{\tm}[1]{\texttt{#1}}
\newcommand{\eps}{\varepsilon}
\newcommand{\First}{\mathrm{First}}
\newcommand{\Follow}{\mathrm{Follow}}
\newcommand{\Predict}{\mathrm{Predict}}
\newcommand{\HALT}{\mathrm{HALT}}
\renewcommand{\checkmark}{\ding{51}}
\renewcommand{\paragraph}[1]{\subsubsection*{#1}}
\fvset{fontsize=\footnotesize, frame=single, framesep=5pt, rulecolor=\color{black!60}}
\setlist{itemsep=2pt, topsep=2pt, parsep=0pt}

% \input primitivo: necesario para leer filas de tabla desde un archivo
\makeatletter
\newcommand{\tabla}[1]{\@@input #1 }
\makeatother

% Generado por scripts/generar_documentacion.py
\newcommand{\NAFN}{319}
\newcommand{\NAFD}{95}
\newcommand{\NMin}{42}
\newcommand{\NFinales}{39}
\newcommand{\NClases}{30}
\newcommand{\NTokens}{26}
\newcommand{\NCeldasDelta}{1260}
\newcommand{\AgrupaIdent}{21}
\newcommand{\AgrupaComent}{30}
\newcommand{\NProd}{43}
\newcommand{\NNoTerm}{20}
\newcommand{\NTerm}{25}
\newcommand{\NCeldasLL}{93}
\newcommand{\NAnulables}{7}
\newcommand{\NConAlternativas}{10}
\newcommand{\TrazaPasos}{21}
\newcommand{\FactorialSlots}{f}
\newcommand{\TrazaVMPasos}{29}
\newcommand{\FibSlots}{t, a}
\newcommand{\FibNodos}{5}
\newcommand{\FibAristas}{9}
\newcommand{\PGOtrosNodos}{31}
\newcommand{\PGOtrasAristas}{14}
\newcommand{\PGSpillsChaitin}{v0}
\newcommand{\PetersenExplorados}{10}
\newcommand{\PendLexer}{0.98}
\newcommand{\RdosLexer}{0.9996}
\newcommand{\PendParser}{1.03}
\newcommand{\RdosParser}{0.9994}
\newcommand{\PendOptTipico}{1.13}
\newcommand{\RdosOptTipico}{0.9631}
\newcommand{\PendOptPeor}{1.81}
\newcommand{\RdosOptPeor}{0.9982}
\newcommand{\PendOptTrabajo}{1.03}
\newcommand{\RdosOptTrabajo}{0.9617}
\newcommand{\PendParserPasos}{1.00}
\newcommand{\RdosParserPasos}{1.0000}
\newcommand{\PasosPorToken}{3.10}
\newcommand{\ColoreoBase}{1.137}
\newcommand{\ColoreoDuplica}{5.4}
\newcommand{\ColoreoRdos}{0.963}
\newcommand{\ProfMaxTokens}{146\,187}
\newcommand{\ProfMaxRI}{44\,873}
\newcommand{\ProfPeorMaxAristas}{1\,124\,250}
\newcommand{\NsPorOpMin}{1048}
\newcommand{\NsPorOpMax}{1602}
\newcommand{\OptVUnoMax}{14.6}
\newcommand{\OptVDosMax}{1.7}
\newcommand{\OptVUnoPend}{1.60}
\newcommand{\NProgValidos}{9}
\newcommand{\NProgErrorlexico}{5}
\newcommand{\NProgErrorsintactico}{8}
\newcommand{\NProgErrorsemantico}{6}
\newcommand{\NProgBucleinfinito}{5}
\newcommand{\NProgramas}{33}
\newcommand{\NTests}{351}


\begin{document}

% =============================================================================
% Portada
\thispagestyle{empty}
\begin{tikzpicture}[remember picture, overlay, x=1pt, y=-1pt, shift={(current page.north west)}]
  \node[anchor=north, inner sep=0] at (306.2,210) {\includegraphics[width=170.5pt]{plantilla/ubp_logo_portada.png}};
  \node[anchor=base, inner sep=0, align=center, text width=432pt, font=\bfseries\fontsize{14}{16.8}\selectfont] at (306.2,459.2) {\Titulo};
  % Un renglón cada 26,6 pt, como en la plantilla; los autores van uno debajo del otro.
  \node[anchor=base west, inner sep=0, font=\fontsize{11}{13.2}\selectfont] at (90,530.1) {%
    \renewcommand{\arraystretch}{1}\setlength{\tabcolsep}{0pt}%
    \begin{tabular}[t]{l}
      \textbf{Asignatura:} \Asignatura \\ \noalign{\vskip 13.4pt}
      \textbf{Actividad:} \Actividad \\ \noalign{\vskip 13.4pt}
      \textbf{Fecha:} \FechaCorta \\ \noalign{\vskip 13.4pt}
      \textbf{\EtiquetaAutor:} \begin{tabular}[t]{@{}l@{}}\Autor\end{tabular} \\ \noalign{\vskip 13.4pt}
      \textbf{\EtiquetaDocente:} \Docente
    \end{tabular}};
\end{tikzpicture}
\null
\newpage

% =============================================================================
% Resumen
\begin{center}
{\bfseries\fontsize{14}{16.8}\selectfont \Titulo\par}
\vspace{22pt}
Por:\par
\vspace{12pt}
\Autor\par
\vspace{18pt}
{\calibri\bfseries\fontsize{12}{14.4}\selectfont RESUMEN\par}
\end{center}
\vspace{-6pt}

El presente trabajo diseña, formaliza, implementa y analiza un compilador y
una máquina virtual abstracta para L-Cero, un lenguaje imperativo mínimo con
variables enteras, condicionales, bucles \tm{while} y operaciones
aritméticas. La implementación se realizó en Python y todos los analizadores
se desarrollaron desde cero, sin generadores automáticos de parsers.

El analizador léxico es un AFD mínimo de \NMin{} estados que se obtiene a
partir de las expresiones regulares de los tokens, y el sintáctico es un
reconocedor LL(1) cuya gramática se demuestra no ambigua con los conjuntos
First y Follow. El código se traduce a una representación intermedia de tres
direcciones y se ejecuta en una máquina de registros con cinta, cabezal,
estados y función de transición explícitos. Ante un programa con errores, el
compilador los informa todos en una sola ejecución, y además emite avisos
sobre construcciones válidas pero sospechosas, como bucles que no terminan.

El trabajo incluye la demostración de que la máquina es Turing-completa, la
reducción desde el problema de la parada que prueba que la terminación de los
programas es indecidible, la demostración de que la asignación de registros
es NP-completa y el contraste entre la complejidad teórica de cada fase y
mediciones reales de tiempo.

\begin{flushright}
\textbf{\textit{\FechaLarga}}
\end{flushright}
\newpage

% =============================================================================
\tableofcontents
\newpage

% =============================================================================
\section{Introducción}

L-Cero es un lenguaje imperativo mínimo que tiene variables enteras,
condicionales \tm{if}/\tm{else}, bucles \tm{while}, operaciones aritméticas
y una instrucción \tm{print}. Este informe presenta el diseño, la
formalización, la implementación y el análisis de un compilador y una
máquina virtual abstracta para ese lenguaje, integrando los cuatro pilares
de la materia: la jerarquía de Chomsky, la teoría de autómatas, la
computabilidad y la complejidad computacional.

La implementación se realizó en Python. Como la consigna prohíbe usar
generadores automáticos de parsers, todos los analizadores se escribieron
desde cero: las expresiones regulares se representan como árboles, y tanto
el AFD del lexer como la tabla LL(1) del parser se calculan con
implementaciones propias de los algoritmos vistos en la materia. Tampoco se
usó el módulo \tm{re} de Python.

El compilador se organiza en las fases de la figura~\ref{fig:fases}. El
código fuente nunca se ejecuta directamente: se traduce a una representación
intermedia (RI) de tres direcciones, el optimizador le asigna los registros
de la máquina y recién entonces se ejecuta en la máquina virtual. El
desarrollo se dividió en tres hitos: el Hito~I abarca las fases 1 a 3
(análisis léxico, sintáctico y semántico), el Hito~II las fases 4 a 6
(representación intermedia, máquina abstracta e indecidibilidad) y el
Hito~III las fases 7 y 8 (complejidad y asignación de registros).

\begin{figure}[H]
\centering
\resizebox{\textwidth}{!}{%
\begin{tikzpicture}[caja/.style={draw, rounded corners=2pt, minimum height=0.9cm, minimum width=2.2cm, align=center, font=\small},
  dato/.style={font=\scriptsize\itshape, align=center}, flecha/.style={-{Stealth[length=2mm]}},
  node distance=0.7cm and 1.1cm]
  \node[caja] (lex) {Lexer\\{\scriptsize AFD mínimo}};
  \node[caja, right=of lex] (par) {Parser\\{\scriptsize LL(1)}};
  \node[caja, right=of par] (sem) {Semántico};
  \node[caja, right=of sem] (ri) {Generador\\de RI};
  \node[caja, right=of ri] (opt) {Optimizador\\{\scriptsize (Chaitin)}};
  \node[caja, right=of opt] (vm) {Máquina\\virtual};
  \draw[flecha] (lex) -- node[dato, above] {tokens} (par);
  \draw[flecha] (par) -- node[dato, above] {AST} (sem);
  \draw[flecha] (sem) -- node[dato, above] {AST} (ri);
  \draw[flecha] (ri) -- node[dato, above] {RI} (opt);
  \draw[flecha] (opt) -- node[dato, above] {RI con\\registros} (vm);
  \node[dato, below=0.1cm of lex] {Tipo 3};
  \node[dato, below=0.1cm of par] {Tipo 2};
  \node[dato, below=0.1cm of sem] {Tipo 1};
  \node[dato, below=0.1cm of opt] {NP-completo};
  \node[dato, below=0.1cm of vm] {Tipo 0};
\end{tikzpicture}}
\caption{Fases del compilador y nivel de la jerarquía de Chomsky de cada una.}
\label{fig:fases}
\end{figure}

% =============================================================================
\section{Desarrollo}

\subsection{El lenguaje L-Cero}

Un programa de ejemplo, que calcula el máximo común divisor por el algoritmo
de Euclides:

\begin{Verbatim}
int a = 1071;
int b = 462;
while (b != 0) {
    int r = a % b;
    a = b;
    b = r;
}
print(a);   // imprime 21
\end{Verbatim}

\paragraph{Alfabeto y tokens}

El alfabeto $\Sigma$ está formado por las letras \tm{a-z}, \tm{A-Z} y
\tm{\_}, los dígitos, los símbolos \tm{+ - * / \% = ! < > ( ) \{ \} ;} y los
blancos (espacio, tabulación, \tm{\textbackslash r} y \tm{\textbackslash n}).
Cualquier otro carácter, por ejemplo \tm{@} o \tm{ñ}, solo puede aparecer
dentro de un comentario; afuera de un comentario es un error léxico.

Para que la tabla de transiciones del lexer no tenga una columna por cada
carácter, los caracteres se agrupan en \NClases{} clases: dos caracteres van
en la misma clase si ninguna expresión regular los distingue. Las letras que
aparecen en las palabras reservadas (\tm{e f h i l n p r s t w}) tienen una
clase cada una, porque el autómata tiene que seguirlas de a una; el resto de
las letras forma una sola clase, lo mismo que los dígitos, y cada símbolo
tiene su propia clase.

La tabla~\ref{tab:er} muestra las expresiones regulares de los tokens. Si
dos tokens reconocen la misma cadena (\tm{while} es una palabra reservada,
pero también cumple la ER de \tm{ID}) gana el de mayor prioridad, y las
palabras reservadas tienen más prioridad que \tm{ID}. Si reconocen cadenas
de distinto largo, gana la más larga (regla del \emph{máximo bocado}):
\tm{whilex} es un identificador y \tm{<=} es un solo token.

\begin{table}[H]
\centering
\begin{grilla}
\begin{tabular}{|L{5.6cm}|L{8.3cm}|}
\hline
\thead{Token} & \thead{Expresión regular} \\ \hline
\tm{INT}, \tm{IF}, \tm{ELSE}, \tm{WHILE}, \tm{PRINT} & la palabra misma: \tm{int}, \tm{if}, \tm{else}, \tm{while}, \tm{print} \\ \hline
\tm{ID} & $\textit{letra}\,(\textit{letra} \mid \textit{dígito})^{*}$ \\ \hline
\tm{NUM} & $\textit{dígito}\;\textit{dígito}^{*}$ \\ \hline
\tm{EQ}, \tm{NEQ}, \tm{LE}, \tm{GE} & \tm{==}, \tm{!=}, \tm{<=}, \tm{>=} \\ \hline
\tm{ASSIGN}, \tm{LT}, \tm{GT} & \tm{=}, \tm{<}, \tm{>} \\ \hline
\tm{PLUS}, \tm{MINUS}, \tm{STAR}, \tm{SLASH}, \tm{PERCENT} & \tm{+}, \tm{-}, \tm{*}, \tm{/}, \tm{\%} \\ \hline
\tm{LPAREN}, \tm{RPAREN}, \tm{LBRACE}, \tm{RBRACE}, \tm{SEMI} & \tm{(}, \tm{)}, \tm{\{}, \tm{\}}, \tm{;} \\ \hline
Comentario (se descarta) & \tm{//}\,$(\Sigma - \{\backslash\tm{n}\})^{*}$ \\ \hline
Espacio (se descarta) & $(\textit{esp} \mid \backslash\tm{n})(\textit{esp} \mid \backslash\tm{n})^{*}$ \\ \hline
\end{tabular}
\end{grilla}
\caption{Expresiones regulares de los \NTokens{} tokens ($\textit{letra} = [\tm{a-zA-Z\_}]$, $\textit{dígito} = [\tm{0-9}]$).}
\label{tab:er}
\end{table}

Otras decisiones del diseño léxico: las palabras reservadas distinguen
mayúsculas (\tm{While} es un identificador); los números son siempre
naturales y los negativos se escriben con el menos unario (\tm{-5} son dos
tokens); no hay operadores lógicos, así que un \tm{!} que no forma \tm{!=}
es un error; y solo hay comentarios de una línea.

\paragraph{Sintaxis y semántica}

La sintaxis del lenguaje, en EBNF, es la siguiente:

\begin{Verbatim}[fontsize=\scriptsize]
programa    = { sentencia } ;
sentencia   = declaracion | asignacion | si | mientras | imprimir | bloque ;
declaracion = "int" id [ "=" expresion ] ";" ;
asignacion  = id "=" expresion ";" ;
si          = "if" "(" expresion ")" bloque [ "else" ( bloque | si ) ] ;
mientras    = "while" "(" expresion ")" bloque ;
imprimir    = "print" "(" expresion ")" ";" ;
bloque      = "{" { sentencia } "}" ;
expresion   = aritmetica [ ( "==" | "!=" | "<" | "<=" | ">" | ">=" ) aritmetica ] ;
aritmetica  = termino { ( "+" | "-" ) termino } ;
termino     = factor { ( "*" | "/" | "%" ) factor } ;
factor      = num | id | "(" expresion ")" | "-" factor ;
\end{Verbatim}

El cuerpo de un \tm{if}, de un \tm{else} y de un \tm{while} tiene que ser
siempre un bloque entre llaves. Con esa decisión no aparece el problema del
\emph{else colgante} y la gramática queda LL(1) sin reglas especiales
(\tm{else if} se permite igual). Los relacionales tienen la menor
precedencia y no se pueden encadenar: \tm{1 < 2 < 3} es un error de
sintaxis.

En cuanto a la semántica, el único tipo es el entero, sin límite de tamaño
(esto es importante para la Turing-completitud de la sección~\ref{sec:maquina});
\tm{/} trunca hacia cero y \tm{\%} da el resto con el signo del dividendo,
como en C; los relacionales valen 1 o 0, y en un \tm{if} o un \tm{while}
cualquier valor distinto de cero cuenta como verdadero. Las reglas sobre
variables se describen en la sección~\ref{sec:semantico}.

% -----------------------------------------------------------------------------
\subsection{Fase 1: análisis léxico}

\paragraph{Construcción del AFD}

El AFD se calcula a partir de las expresiones regulares en tres pasos,
implementados en \tm{lcero/automata.py}:

\begin{enumerate}
  \item \textbf{Construcción de Thompson.} Cada ER se convierte en un AFN-$\eps$
  armando un fragmento por cada operación (símbolo, concatenación, unión y
  clausura de Kleene). Los \NTokens{} fragmentos se unen con un estado
  inicial nuevo y transiciones $\eps$. El resultado es un AFN-$\eps$ de
  \NAFN{} estados, donde cada estado final indica qué token reconoce y con
  qué prioridad.
  \item \textbf{Construcción por subconjuntos.} Cada estado del AFD es un
  conjunto de estados del AFN cerrado por $\eps$:
  $\delta_D(S, a) = \eps\text{-clausura}(\mathrm{mover}(S, a))$. Si un
  conjunto contiene finales de varios tokens, se queda con el de mayor
  prioridad (así \tm{while} da \tm{WHILE} y no \tm{ID}). El conjunto vacío
  funciona como estado sumidero. Resultan \NAFD{} estados.
  \item \textbf{Minimización de Hopcroft.} Se parte de una partición que
  separa los estados no finales y, entre los finales, uno por cada token, y
  se refina hasta que no se puede partir más. Cada bloque de la partición
  final es un estado del AFD mínimo, que tiene \NMin{} estados.
\end{enumerate}

\begin{definicion}[AFD mínimo del lexer]
$M = (Q, C, \delta, q_0, F)$, donde
$Q = \{q_0, \dots, q_{40}, q_\emptyset\}$ ($|Q| = \NMin$ y $q_\emptyset$ es
el sumidero); $C$ es el conjunto de las \NClases{} clases de caracteres;
$\delta : Q \times C \to Q$ es la función de transición, dada por la matriz
del anexo~A; $q_0$ es el estado inicial; y
$F = Q \setminus \{q_0, q_{13}, q_\emptyset\}$ ($|F| = \NFinales$). Cada
estado final tiene asociado el token que acepta.
\end{definicion}

La figura~\ref{fig:afd} muestra la parte del autómata que reconoce
identificadores, palabras reservadas y números. La parte de los operadores
y la matriz $\delta$ completa ($\NMin \times \NClases$) están en el
anexo~A. Por ejemplo, $q_{24}$ representa haber leído \tm{in}: con una
\tm{t} pasa a $q_{33}$ (\tm{INT}) y con cualquier otra letra o dígito pasa
a $q_2$, que es el estado de identificador.

\begin{figure}[H]
\centering
\includegraphics[width=\textwidth]{generado/afd_palabras.pdf}
\caption{Parte del AFD mínimo que reconoce identificadores, palabras reservadas y números.}
\label{fig:afd}
\end{figure}

\paragraph{Analizador léxico}

Para reconocer los tokens, el lexer recorre el AFD desde la posición actual
y recuerda el último estado final por el que pasó. Cuando el autómata cae en
el sumidero (o se termina la entrada) emite el token de ese último estado
final y continúa desde donde terminó. Los comentarios y los espacios se
reconocen pero no se pasan al parser. Si no se llega a ningún estado final,
hay un error léxico, que se informa con la línea y la columna; por ejemplo:
\emph{«el carácter '@' no pertenece al alfabeto de L-Cero»}.

Para verificar el autómata, las pruebas lo comparan con otros métodos: el
AFD mínimo acepta lo mismo que el de subconjuntos, el algoritmo de Moore no
encuentra estados equivalentes (es decir, es realmente mínimo) y el lexer da
los mismos tokens que simular directamente el AFN-$\eps$ sobre 2000 cadenas
al azar.

% -----------------------------------------------------------------------------
\subsection{Fase 2: análisis sintáctico}

\paragraph{Gramática}

Para pasar la EBNF a una gramática LL(1) se reemplazaron las repeticiones
$\{\ \}$ por no terminales recursivos a derecha (\nt{ListaSent},
\nt{AritCola}, \nt{TermCola}) y los opcionales $[\ ]$ por no terminales con
una alternativa $\eps$ (\nt{DeclInit}, \nt{SiNo}, \nt{ExprRel}). Así no
queda recursión a izquierda. La gramática tiene \NNoTerm{} no terminales,
\NTerm{} terminales (contando \tm{\$}, el fin de entrada) y \NProd{}
producciones:

\begin{center}
\small
\begin{tabular}{@{}l c p{11.6cm}@{}}
\tabla{generado/f2_gramatica_compacta.tex}
\end{tabular}
\end{center}

\paragraph{Conjuntos First y Follow}

$\First(\alpha)$ es el conjunto de terminales con los que puede empezar algo
derivado de $\alpha$ (más $\eps$ si $\alpha$ puede derivar la cadena vacía),
y $\Follow(A)$ son los terminales que pueden aparecer inmediatamente después
de $A$ en alguna forma sentencial. Se calculan aplicando las siguientes
reglas hasta que ningún conjunto cambia:

\begin{itemize}
  \item $\First(A)$ es la unión de $\First(\alpha)$ para cada producción
  $A \to \alpha$, donde $\First(X_1 \cdots X_m)$ incluye $\First(X_1)$ y, si
  $X_1$ es anulable, también $\First(X_2 \cdots X_m)$, y así sucesivamente.
  \item $\tm{\$} \in \Follow(\nt{Programa})$; si $A \to \alpha B \beta$,
  entonces $\First(\beta) - \{\eps\} \subseteq \Follow(B)$; y si además
  $\beta$ es anulable, $\Follow(A) \subseteq \Follow(B)$.
\end{itemize}

La tabla~\ref{tab:ff} muestra el resultado. El cálculo paso a paso, con lo
que aporta cada producción, está en el anexo~B.

\begin{table}[H]
\centering
\begin{grilla}
\begin{tabular}{|l|C{1.1cm}|L{4.8cm}|L{5.5cm}|}
\hline
\thead{$A$} & \thead{Anul.} & \thead{$\First(A)$} & \thead{$\Follow(A)$} \\ \hline
\tabla{generado/f2_first_follow_grid.tex}
\end{tabular}
\end{grilla}
\caption{Anulables, First y Follow de cada no terminal.}
\label{tab:ff}
\end{table}

\paragraph{La gramática no es ambigua}

Para cada producción se define
$\Predict(A \to \alpha) = \First(\alpha) - \{\eps\}$, más $\Follow(A)$ si
$\alpha$ es anulable. La gramática es LL(1) si las producciones de un mismo
no terminal tienen conjuntos Predict disjuntos. La tabla~\ref{tab:predict}
muestra los Predict de los \NConAlternativas{} no terminales que tienen más
de una producción, y en todos los casos son disjuntos. El caso más
importante es \nt{SiNo}: como los cuerpos del \tm{if} terminan siempre en
\tm{\}}, \tm{else} no pertenece a $\Follow(\nt{SiNo})$ y no hay conflicto.
Si se permitieran cuerpos sin llaves, \tm{else} sí estaría en ese Follow y
aparecería el conflicto del else colgante (una de las pruebas lo verifica
con la gramática ambigua).

\begin{table}[H]
\centering
\begin{grilla}
\begin{tabular}{|l|L{11.1cm}|}
\hline
\thead{$A$} & \thead{Predict de cada producción de $A$ (con su número)} \\ \hline
\tabla{generado/f2_predict_compacto_grid.tex}
\end{tabular}
\end{grilla}
\caption{Conjuntos Predict de los no terminales con más de una producción.}
\label{tab:predict}
\end{table}

\begin{teorema}
Una gramática LL(1) no es ambigua. En particular, la gramática de L-Cero no
es ambigua.
\end{teorema}

\begin{proof}
Supóngase que una cadena $w$ tiene dos derivaciones por izquierda distintas
y considérese el primer paso en que se separan. Ambas llegaron a la misma
forma sentencial $xA\gamma$ (con $x$ formado por terminales y $w = xy$); una
aplica $A \to \alpha$ y la otra $A \to \beta$. Sea $a$ el primer símbolo de
$y\tm{\$}$. Si $a$ proviene de $\alpha$, entonces $a \in \First(\alpha)$; si
no, $\alpha$ es anulable y $a$ proviene de lo que sigue a $A$, por lo que
$a \in \Follow(A)$. En ambos casos $a \in \Predict(A \to \alpha)$, y por el
mismo razonamiento $a \in \Predict(A \to \beta)$, lo que contradice que los
Predict sean disjuntos. Como la tabla~\ref{tab:predict} muestra que la
gramática de L-Cero es LL(1), no es ambigua.
\end{proof}

Con los Predict se arma la tabla LL(1) $M[A, a]$ (anexo~C), que tiene
\NCeldasLL{} celdas ocupadas y ninguna con dos producciones.

\paragraph{Reconocedor sintáctico}

El parser (\tm{lcero/parser.py}) es un autómata de pila con un solo estado,
$P = (\{q\}, T, N \cup T \cup \{\tm{\$}\}, \delta, q, \tm{\$}, \emptyset)$,
que acepta por pila vacía. Si en el tope hay un no terminal $A$ y el
próximo token es $a$, reemplaza $A$ por el lado derecho de $M[A, a]$
(invertido, para que el primer símbolo quede arriba); si en el tope hay un
terminal igual al próximo token, lo desapila y avanza. Como decide mirando
un solo token, es determinista. La tabla~\ref{tab:traza} muestra la traza
para \texttt{x = a * 2;} (el tope de la pila está a la
derecha). Mientras reconoce, el parser construye el árbol de derivación, que
luego se transforma en un árbol de sintaxis abstracta (AST).

\begin{table}[H]
\centering
\begin{grilla}\scriptsize\renewcommand{\arraystretch}{1.3}
\begin{tabular}{|r|l|r|l|}
\hline
\thead{Paso} & \thead{Pila} & \thead{Entrada} & \thead{Acción} \\ \hline
\tabla{generado/f2_traza_grid.tex}
\end{tabular}
\end{grilla}
\caption{Traza del autómata de pila.}
\label{tab:traza}
\end{table}

Cuando la celda $M[A, a]$ está vacía hay un error de sintaxis. Para que el
mensaje sea útil, si $A$ es anulable se aplica igual $A \to \eps$ y el error
se detecta más abajo en la pila, sin consumir el token. Por ejemplo, si
falta un punto y coma el mensaje es \emph{«se esperaba un operador o ';',
pero se encontró 'print'»}, con la línea y la columna.

\paragraph{Recuperación de errores}

Para no detenerse en el primer error, el parser usa recuperación en modo
pánico, que se apoya en los conjuntos First y Follow calculados antes:

\begin{itemize}
  \item Si en el tope hay un terminal distinto del próximo token, se supone que
  ese terminal faltaba y se lo desapila, como si se hubiera insertado.
  \item Si la celda $M[A, a]$ está vacía, se descartan tokens hasta encontrar
  uno de $\First(A)$, y entonces se vuelve a intentar $A$, o uno de
  $\Follow(A)$, y entonces se abandona $A$. El fin de entrada nunca se
  descarta.
\end{itemize}

Cada paso consume un token o achica la pila, así que el análisis siempre
termina. Para no informar errores en cascada, después de cada error no se
informan otros hasta emparejar tres tokens (la misma regla que usa Bison).
De esta forma, un programa con varios errores los muestra todos en una sola
ejecución, como en la figura~\ref{fig:errores}.

\begin{figure}[H]
\begin{Verbatim}[breaklines]
int x = 1
int y = x +;
print(y);
while (x < 3 {
    x = x + 1;
}
\end{Verbatim}
\vspace{2pt}
\begin{Verbatim}[breaklines]
Error sintáctico (línea 2, columna 1): se esperaba un operador o ';', pero se encontró 'int'
Error sintáctico (línea 2, columna 12): se esperaba una expresión, pero se encontró ';'
Error sintáctico (línea 4, columna 14): se esperaba un operador aritmético o ')', pero se encontró '{'
3 errores
\end{Verbatim}
\caption{Programa con tres errores de sintaxis (arriba) y lo que informa el compilador (abajo).}
\label{fig:errores}
\end{figure}

% -----------------------------------------------------------------------------
\subsection{Fase 3: análisis semántico}
\label{sec:semantico}

El análisis semántico (\tm{lcero/semantic.py}) recorre el AST y controla
tres reglas: toda variable debe estar declarada antes de usarse; no se puede
declarar dos veces una variable visible, tampoco en un bloque anidado (como
en Java); y no se puede leer una variable que, en algún camino posible,
todavía no tiene valor. Los ámbitos se manejan con un diccionario de nombres
visibles y una pila con lo declarado en cada bloque. Para la inicialización
se lleva el conjunto de variables que seguro tienen valor: después de un
\tm{if} con \tm{else} es la intersección de lo que dejan las dos ramas, y
después de un \tm{if} sin \tm{else} o de un \tm{while} es el mismo que
antes, porque el cuerpo podría no ejecutarse.

Esta fase tampoco se detiene en el primer error: cada error se registra y el
análisis sigue con el resto del programa. Para no repetir el mismo problema,
cada variable sin declarar y cada variable sin inicializar se informa una
sola vez, y una doble declaración no reemplaza a la declaración original.

\paragraph{Por qué esta fase excede a un autómata con pila}

Sea $L_{\mathrm{sem}}$ el conjunto de los programas que pasan las tres
primeras fases.

\begin{teorema}
$L_{\mathrm{sem}}$ no es libre de contexto.
\end{teorema}

\begin{proof}
Sea el lenguaje regular
$R = \tm{int}\,\sqcup\,(\tm{a}|\tm{b})^{+}\tm{;}(\tm{a}|\tm{b})^{+}\tm{=0;}$,
donde $\sqcup$ es un espacio. Un programa de la forma
$\tm{int}\sqcup u\tm{;}v\tm{=0;}$ es válido solo si $v$ está declarada, es
decir, si $v = u$. Entonces
$L' = L_{\mathrm{sem}} \cap R = \{\tm{int}\sqcup w\tm{;}w\tm{=0;} : w \in \{\tm{a},\tm{b}\}^{+}\}$.
Si $L_{\mathrm{sem}}$ fuera libre de contexto, $L'$ también lo sería
(intersección con un regular). Se aplica el lema de bombeo con
$s = \tm{int}\sqcup\tm{a}^p\tm{b}^p\tm{;}\tm{a}^p\tm{b}^p\tm{=0;}$, donde $p$
es la constante del lema, y una descomposición $s = uvxyz$ con
$|vxy| \le p$ y $|vy| \ge 1$:
\begin{itemize}
  \item si $v$ o $y$ contienen algún símbolo fijo (\tm{i}, \tm{n}, \tm{t},
  $\sqcup$, \tm{;}, \tm{=}, \tm{0}), al bombear con $i = 0$ la cadena deja de
  pertenecer a $R$;
  \item si $vxy$ está dentro de una sola copia de $w$, al bombear cambia esa
  copia y la otra no, por lo que las copias quedan distintas;
  \item si $vxy$ abarca el final de la primera copia y el principio de la
  segunda, $v$ está entre las \tm{b} de la primera e $y$ entre las \tm{a} de
  la segunda. Con $i = 2$, la primera copia tiene más \tm{b} que la segunda
  (si $|v| > 0$) o la segunda tiene más \tm{a} (si $|v| = 0$).
\end{itemize}
En todos los casos $uv^2xy^2z$ o $uv^0xy^0z$ no pertenece a $L'$, así que
$L'$ no es libre de contexto, y $L_{\mathrm{sem}}$ tampoco.
\end{proof}

En consecuencia, ningún autómata de pila (ni ningún parser LL o LR) puede
controlar que las variables estén declaradas: hay que comparar dos nombres
que pueden estar arbitrariamente lejos, y una pila no lo permite sin perder
lo que tiene debajo. Sí alcanza un autómata linealmente acotado: cada
control se puede hacer recorriendo la entrada y marcando celdas, sin usar
más espacio que el de la entrada (comparar cada uso con las declaraciones
anteriores letra por letra, contar llaves para saber si la declaración es
visible y buscar caminos en el grafo de flujo para la inicialización). Por
el teorema de Kuroda, $L_{\mathrm{sem}}$ es sensible al contexto (Tipo~1).

\paragraph{Avisos del compilador}

Además de los errores, el compilador emite avisos sobre construcciones
válidas pero casi seguramente equivocadas (\tm{lcero/lint.py}). Los avisos
no rechazan el programa: se muestran al ejecutarlo (la opción
\tm{--sin-avisos} los oculta) y con el comando \tm{lcero avisos}. Se
detectan cuatro casos:

\begin{itemize}
  \item variables que se declaran y nunca se usan, o que reciben valores que
  nunca se leen;
  \item condiciones constantes en un \tm{if} o un \tm{while}, es decir, una
  rama que nunca se ejecuta o un bucle que no termina;
  \item condiciones de \tm{while} en las que ninguna variable se asigna dentro
  del cuerpo: como L-Cero no tiene \tm{break} ni expresiones con efectos, la
  condición vale siempre lo mismo y, si el bucle entra, no termina;
  \item código inalcanzable, es decir, sentencias que están después de una que
  nunca termina.
\end{itemize}

El programa de la figura~\ref{fig:avisos} produce los cuatro tipos de aviso.

\begin{figure}[H]
\begin{Verbatim}[breaklines]
int i = 0;
int suma = 0;
int sinUso = 5;
while (i < 10) {
    suma = suma + 2;
}
if (1 > 2) {
    print(suma);
}
while (1 == 1) {
}
print(i);
\end{Verbatim}
\vspace{2pt}
\begin{Verbatim}[breaklines]
Aviso (línea 3, columna 5): la variable 'sinUso' recibe valores que nunca se leen
Aviso (línea 4, columna 10): ninguna variable de la condición cambia en el cuerpo: si el bucle entra, no termina nunca
Aviso (línea 7, columna 7): la condición del if es siempre falsa: el bloque nunca se ejecuta
Aviso (línea 10, columna 10): la condición del while es siempre verdadera: el bucle no termina nunca
Aviso (línea 12, columna 1): código inalcanzable: la sentencia anterior no termina nunca
\end{Verbatim}
\caption{Programa con los cuatro tipos de aviso (arriba) y los avisos del compilador (abajo).}
\label{fig:avisos}
\end{figure}

En la sección~\ref{sec:indecidibilidad} se explica por qué estos avisos no
pueden detectar todos los bucles infinitos.

% -----------------------------------------------------------------------------
\subsection{Fase 4: representación intermedia}

El compilador traduce el AST a código de tres direcciones (\tm{lcero/ir.py}).
Los operandos son variables, temporales (\tm{\$t1}, \tm{\$t2}, \dots) o
constantes, y en esta etapa no hay límite de registros. La
tabla~\ref{tab:ri} resume las instrucciones.

\begin{table}[H]
\centering
\begin{grilla}
\begin{tabular}{|L{5.4cm}|L{8.5cm}|}
\hline
\thead{Instrucción} & \thead{Significado} \\ \hline
\tm{x = a op b} & Operación binaria (aritmética o relacional) \\ \hline
\tm{x = -a}, \tm{x = a} & Menos unario y copia \\ \hline
\tm{L:}, \tm{goto L}, \tm{ifz a goto L} & Etiqueta, salto y salto si $a = 0$ \\ \hline
\tm{print a}, \tm{halt} & Salida y fin del programa \\ \hline
\tm{x = load [n]}, \tm{store [n] = a} & Lectura y escritura de la cinta (solo las agrega el optimizador) \\ \hline
\end{tabular}
\end{grilla}
\caption{Instrucciones de la representación intermedia.}
\label{tab:ri}
\end{table}

Un \tm{while} se traduce como \tm{L1:} condición; \tm{ifz t goto L2};
cuerpo; \tm{goto L1}; \tm{L2:}, y un \tm{if} de forma similar. Durante la
generación se pliegan las constantes (\tm{2 * (3 + 4)} se traduce como
\tm{14}). La figura~\ref{fig:ri} muestra el resultado para el factorial.

\begin{figure}[H]
\centering
\begin{minipage}[t]{0.42\textwidth}
\VerbatimInput[firstline=3]{../../tests/programas/validos/factorial.l0}
\end{minipage}\hspace{1cm}
\begin{minipage}[t]{0.42\textwidth}
\VerbatimInput{generado/factorial_ri.txt}
\end{minipage}
\caption{Programa del factorial (izquierda) y su RI (derecha).}
\label{fig:ri}
\end{figure}

Esta RI todavía usa nombres ilimitados. El optimizador de la
sección~\ref{sec:registros} le asigna los registros de la máquina, y eso es
lo que finalmente se ejecuta.

% -----------------------------------------------------------------------------
\subsection{Fase 5: máquina virtual}
\label{sec:maquina}

La máquina (\tm{lcero/vm.py}) es una máquina de registros con una cinta,
modelada con los mismos elementos que una máquina de Turing.

\begin{definicion}[Máquina de L-Cero]
Para un programa de $m$ instrucciones y $k \ge 2$ registros, la máquina es
$\mathcal{M} = (Q, k, \Gamma, \delta, q_0, q_f)$, donde
$Q = \{q_0, \dots, q_{m-1}, q_f\}$ tiene un estado por instrucción más el
estado de parada; hay $k$ registros $R_0, \dots, R_{k-1}$ con enteros sin
cota; la cinta tiene celdas numeradas $0, 1, 2, \dots$ con símbolos de
$\Gamma = \mathbb{Z}$ (el blanco es 0) y un cabezal que se mueve de a una
celda; y $\delta$ es la función de transición de la tabla~\ref{tab:delta-vm}.
Una configuración es $C = (q, \rho, h, \tau, \omega)$: estado, contenido de
los registros, posición del cabezal, contenido de la cinta y valores
impresos hasta el momento.
\end{definicion}

\begin{table}[H]
\centering
\begin{grilla}
\begin{tabular}{|L{4.2cm}|L{9.7cm}|}
\hline
\thead{Instrucción en $q$} & \thead{Nueva configuración} \\ \hline
\tm{mov Rd, A} & $(q{+}1,\ \rho[d \mapsto A],\ h,\ \tau,\ \omega)$ \\ \hline
\tm{add/sub/mul/\dots{} Rd, A, B} & $(q{+}1,\ \rho[d \mapsto A \mathbin{op} B],\ h,\ \tau,\ \omega)$; error si se divide por 0 \\ \hline
\tm{load Rd, [a]} & si $h < a$: $h{+}1$; si $h > a$: $h{-}1$; si $h = a$: $(q{+}1,\ \rho[d \mapsto \tau(h)],\ h,\ \tau,\ \omega)$ \\ \hline
\tm{store [a], A} & mueve el cabezal igual que \tm{load}; si $h = a$: $(q{+}1,\ \rho,\ h,\ \tau[h \mapsto A],\ \omega)$ \\ \hline
\tm{jmp q'} / \tm{jz A, q'} & pasa a $q'$ (en \tm{jz}, solo si $A = 0$; si no, a $q{+}1$) \\ \hline
\tm{print A} & $(q{+}1,\ \rho,\ h,\ \tau,\ \omega \cdot A)$ \\ \hline
\tm{halt} & pasa a $q_f$ \\ \hline
\end{tabular}
\end{grilla}
\caption{Función de transición de la máquina ($A$ y $B$ son registros o constantes).}
\label{tab:delta-vm}
\end{table}

Igual que en una MT, para leer o escribir la celda $a$ el cabezal tiene que
llegar hasta ella moviéndose una celda por paso. En la clase \tm{Machine},
el método \tm{delta} aplica exactamente un paso de la tabla y \tm{run} lo
repite hasta llegar a $q_f$ (o hasta un límite de pasos). Por ejemplo, el
programa \tm{int a = 7; int b = 5; int c = 2; print(a * b - c); \dots} con
solo 2 registros obliga a guardar las variables en la cinta, y en la traza
de la figura~\ref{fig:traza-vm} se ve el movimiento del cabezal.

\begin{figure}[H]
\centering
\begin{minipage}[t]{0.33\textwidth}
\VerbatimInput[fontsize=\scriptsize,lastline=12]{generado/traza_maquina.txt}
\end{minipage}\hfill
\begin{minipage}[t]{0.63\textwidth}
\VerbatimInput[fontsize=\scriptsize,lastline=12]{generado/traza_configuraciones.txt}
\end{minipage}
\caption{Programa de la máquina (izquierda) y primeras configuraciones de la ejecución (derecha).}
\label{fig:traza-vm}
\end{figure}

\paragraph{Turing-completitud}

\begin{teorema}
La máquina de L-Cero es Turing-completa para cualquier $k \ge 2$.
\end{teorema}

\begin{proof}
Minsky demostró que las máquinas de dos contadores son Turing-completas.
Esas máquinas tienen dos registros naturales y tres instrucciones:
$\mathrm{INC}(r, j)$, que suma 1 a $r$ y salta a $j$;
$\mathrm{DEC}(r, j, \ell)$, que salta a $\ell$ si $r = 0$ y, si no, le resta
1 y salta a $j$; y $\mathrm{HALT}$. Cada una se traduce a la máquina de
L-Cero con a lo sumo tres instrucciones: \tm{add Rr, Rr, 1; jmp} para INC,
\tm{jz Rr; sub Rr, Rr, 1; jmp} para DEC y \tm{halt}. Como los registros no
tienen cota, la simulación es exacta paso a paso, así que la máquina puede
simular cualquier máquina de dos contadores y, por lo tanto, cualquier MT.
\end{proof}

Para la Fase~6 hace falta algo más fuerte: que el propio lenguaje L-Cero
pueda simular cualquier MT. En \tm{lcero/turing.py} se implementó una
función $f$ que, dada una MT $M$ y una entrada $w$, genera un programa
L-Cero que la simula. La cinta de $M$ se guarda en tres variables: \tm{cur}
(el símbolo bajo el cabezal), \tm{izq} y \tm{der} (lo que hay a cada lado,
escrito como un número en base $b = |\Gamma|$, con la celda más cercana al
cabezal como dígito menos significativo). Mover el cabezal a la derecha es
\tm{izq = izq * b + cur; cur = der \% b; der = der / b}, y a la izquierda es
lo simétrico. El programa es un \tm{while (estado != H)} con un \tm{if} por
cada estado y cada símbolo (un ejemplo generado es \tm{maquina\_turing.l0},
en las pruebas). Por inducción en la cantidad de pasos se ve que, después
de $t$ vueltas del bucle, las variables representan la configuración de $M$
después de $t$ pasos; por lo tanto, \textbf{el programa termina si y solo si
$M$ se detiene con $w$}. Esto se probó con los castores afanosos de 2 y 3
estados (que dan 4 unos en 6 pasos y 6 unos en 14 pasos, igual que la MT) y
con 60 máquinas al azar.

Todo esto depende de que los enteros no tengan cota. Con enteros de 32 bits
la máquina tendría una cantidad finita de configuraciones, sería en el fondo
un autómata finito y se podría decidir si termina.

% -----------------------------------------------------------------------------
\subsection{Fase 6: indecidibilidad de la terminación}
\label{sec:indecidibilidad}

La consigna de esta fase es: \emph{«Demuestre formalmente que es imposible
construir un algoritmo estático global capaz de determinar si cualquier
programa dado escrito en L-Cero detendrá su ejecución en la Máquina Virtual
sin caer en un bucle infinito»}.

Sean $\HALT_{MT} = \{\langle M, w\rangle : M$ se detiene con entrada $w\}$
y $\HALT_{L0} = \{P : P$ es un programa L-Cero válido cuya ejecución en la
máquina termina$\}$. Un algoritmo como el de la consigna sería un decisor
para $\HALT_{L0}$.

\begin{teorema}
$\HALT_{L0}$ es indecidible.
\end{teorema}

\begin{proof}
Se reduce el problema de la parada, que Turing probó indecidible, a
$\HALT_{L0}$, usando la función $f$ de la sección anterior:

\begin{itemize}
  \item $f$ es computable y total: para cualquier $\langle M, w\rangle$
  escribe un programa con un bloque por cada par (estado, símbolo) de $M$ y
  codifica $w$ como un número. Es un procedimiento mecánico que siempre
  termina (la función \tm{to\_lcero}).
  \item $f(\langle M, w\rangle)$ es un programa L-Cero válido: declara todas
  sus variables inicializadas y no puede dar errores de ejecución, porque
  solo divide por la constante $b \ge 1$.
  \item $\langle M, w\rangle \in \HALT_{MT} \iff f(\langle M, w\rangle) \in \HALT_{L0}$,
  porque el programa termina si y solo si $M$ se detiene con $w$.
\end{itemize}

Supóngase que existe un algoritmo $D$ que decide $\HALT_{L0}$. Entonces,
para decidir si $M$ se detiene con $w$ alcanzaría con calcular
$f(\langle M, w\rangle)$ y preguntarle a $D$. Eso decidiría el problema de
la parada, lo cual es imposible. Por lo tanto $D$ no existe: no hay ningún
algoritmo, estático ni de otro tipo, que decida para todo programa L-Cero si
su ejecución termina.
\end{proof}

Con la misma idea se prueba que tampoco se puede decidir si un programa
imprime algo (basta agregar un \tm{print} al final de
$f(\langle M, w\rangle)$); en general, es el teorema de Rice. Lo que sí se
puede es \emph{semidecidir} la terminación: si al ejecutar el programa este
termina, se sabe, pero si no termina nunca se tiene certeza. Por eso
\tm{Machine.run} acepta un límite de pasos, y en las pruebas de bucles
infinitos solo se puede verificar que no terminaron en 200\,000 pasos.

El teorema también explica el alcance de los avisos de la Fase~3. Los avisos
sobre bucles que no terminan son seguros pero no completos: si el compilador
avisa que un bucle no termina, es cierto (en el caso de la condición
invariante, siempre que el bucle llegue a entrar), pero no detecta todos los
bucles infinitos. De los cinco programas de \tm{tests/programas/bucle\_infinito}
solo avisa en \tm{condicion\_constante.l0}. Por ejemplo, en
\tm{olvido\_de\_incremento.l0} el contador \tm{i} solo se incrementa dentro
de un \tm{if} cuya condición nunca se cumple; como la variable sí se asigna en
el cuerpo del bucle, el análisis no puede descartar que el bucle termine y no
da ningún aviso. Un análisis que detectara todos los casos decidiría
$\HALT_{L0}$, lo cual es imposible.

% -----------------------------------------------------------------------------
\subsection{Fase 7: análisis asintótico y profiling}

Se llama $n$ al largo de la entrada en caracteres, $t$ a la cantidad de
tokens, $N$ a la cantidad de instrucciones de la RI, $V$ a la cantidad de
variables y temporales, $E$ a las aristas del grafo de interferencia y $L$
a la máxima cantidad de variables vivas al mismo tiempo. La
tabla~\ref{tab:complejidad} resume la complejidad de cada componente.

\begin{table}[H]
\centering
\begin{grilla}
\begin{tabular}{|L{3.1cm}|L{3.1cm}|L{2.1cm}|L{2.6cm}|L{2.3cm}|}
\hline
\thead{Componente} & \thead{$O$} & \thead{$\Omega$} & \thead{$\Theta$} & \thead{Espacio} \\ \hline
Lexer & $O(n)$ & $\Omega(n)$ & $\Theta(n)$ & $\Theta(n)$ \\ \hline
Parser & $O(t)$ & $\Omega(t)$ & $\Theta(t)$ & $\Theta(t)$ \\ \hline
Optimizador (caso general) & $O(N V L + (V{+}E)\log V)$ por ronda & $\Omega(N + V)$ & Depende del programa & $\Theta(N L + V + E)$ \\ \hline
Optimizador, pocas variables vivas & $O(N \log N)$ & $\Omega(N)$ & Lineal en la práctica & $\Theta(N)$ \\ \hline
Optimizador, todas vivas ($K_m$) & $O(N^2 \log N)$ & $\Omega(N^2)$ & Domina $N^2$ & $\Theta(N^2)$ \\ \hline
\end{tabular}
\end{grilla}
\caption{Complejidad temporal y espacial de cada componente.}
\label{tab:complejidad}
\end{table}

\paragraph{Lexer}

En general, el máximo bocado puede tener que retroceder varios caracteres.
En este AFD eso no ocurre: los únicos estados no finales son $q_0$,
$q_{13}$ (el \tm{!}) y el sumidero, y a $q_{13}$ solo se llega desde $q_0$.
Entonces, una vez que se pasa por un estado final, todos los siguientes
también son finales hasta caer en el sumidero, y lo único que se vuelve a
leer es el carácter que provocó la caída. Cada carácter se lee a lo sumo dos
veces, así que el tiempo es $O(n)$, y como hay que leerlos todos también es
$\Omega(n)$.

\paragraph{Parser}

El autómata hace un paso por cada nodo del árbol de derivación. La gramática
no tiene ciclos ($A \Rightarrow^{+} A$) y cada producción tiene a lo sumo un
símbolo anulable acompañado de otro que no lo es, así que el árbol tiene
$O(t)$ nodos. Cada token se consume una sola vez, lo que da $\Theta(t)$. En
las mediciones se obtuvieron \PasosPorToken{} pasos por token en promedio.

\paragraph{Optimizador}

El análisis de vivacidad visita cada instrucción pocas veces (se midieron
alrededor de $1.3\,N$ visitas) con conjuntos de a lo sumo $L$ variables;
armar el grafo cuesta $O(N \cdot L)$; y colorearlo con el montículo de
candidatos cuesta $O((V+E)\log V)$. Si hay pocas variables vivas a la vez,
el grafo tiene $O(N)$ aristas y todo es casi lineal. En el peor caso, si
$m$ variables están vivas al mismo tiempo, el grafo es completo y tiene
$m(m-1)/2$ aristas, por lo que el costo es al menos cuadrático.

\paragraph{Profiling}

Para contrastar con la teoría se generaron programas L-Cero válidos de
distintos tamaños (\tm{scripts/profiling.py}) y se midió el tiempo de cada
fase (mediana de 5 ejecuciones). Se usaron programas típicos de hasta
\ProfMaxTokens{} tokens y, para el peor caso del optimizador, programas con
todas las variables vivas a la vez. La figura~\ref{fig:profiling} muestra
las mediciones junto con la curva teórica ajustada, y la
tabla~\ref{tab:pendientes} la pendiente en escala log-log, que estima el
exponente (1 es lineal y 2 es cuadrático).

\begin{figure}[H]
\centering
\includegraphics[width=0.98\textwidth]{generado/profiling_fases.pdf}
\caption{Tiempo medido (puntos) contra el modelo teórico (línea).}
\label{fig:profiling}
\end{figure}

\begin{table}[H]
\centering
\begin{grilla}
\begin{tabular}{|l|c|c|c|}
\hline
\thead{Serie} & \thead{Modelo} & \thead{Pendiente log-log} & \thead{$R^2$} \\ \hline
Lexer & $c \cdot n$ & \PendLexer & \RdosLexer \\ \hline
Parser & $c \cdot t$ & \PendParser & \RdosParser \\ \hline
Optimizador, programas típicos & $c \cdot N$ & \PendOptTipico & \RdosOptTipico \\ \hline
Optimizador, peor caso & $c \cdot N^2$ & \PendOptPeor & \RdosOptPeor \\ \hline
\end{tabular}
\end{grilla}
\caption{Ajuste de las mediciones al modelo teórico.}
\label{tab:pendientes}
\end{table}

El lexer y el parser dan prácticamente 1, como se esperaba. En el
optimizador con programas típicos la pendiente es \PendOptTipico: si en
lugar del tiempo se cuentan las operaciones (visitas, nodos y aristas), da
\PendOptTrabajo, así que el algoritmo es lineal y la diferencia se debe a
que el tiempo por operación crece un poco con el tamaño (efecto de la
memoria en Python). En el peor caso da \PendOptPeor, cerca de 2.

El profiling también permitió encontrar un error. En la primera versión del
optimizador, cada vez que había que elegir una variable para mandar a
memoria se recorrían todos los nodos buscando la más barata. Se había
considerado un costo menor, pero en programas largos se eligen muchas
variables y el algoritmo resultaba cuadrático: pendiente \OptVUnoPend{} y
\OptVUnoMax{} segundos para el programa más grande. Al reemplazar esa
búsqueda por un montículo el tiempo bajó a \OptVDosMax{} segundos, con
exactamente las mismas decisiones.

% -----------------------------------------------------------------------------
\subsection{Fase 8: asignación de registros}
\label{sec:registros}

La RI usa infinitos nombres, pero la máquina tiene $k$ registros. Una
variable está \emph{viva} en un punto si su valor actual se va a usar más
adelante. Dos variables \emph{interfieren} si una se define en un punto
donde la otra está viva; en ese caso no pueden compartir registro. Con esa
relación se arma el \emph{grafo de interferencia} $G(P)$, y asignar
registros sin mandar nada a memoria equivale a colorear $G(P)$ con $k$
colores.

\begin{teorema}
Para $k \ge 3$, decidir si las variables de un programa entran en $k$
registros (sin mandar ninguna a memoria) es NP-completo.
\end{teorema}

\begin{proof}
\emph{Pertenece a NP}: dada una asignación de registros, se calcula $G(P)$
en tiempo polinomial y se verifica que ninguna arista una dos variables con
el mismo registro.

\emph{Es NP-difícil}: se reduce $k$-coloreo, que es NP-completo para
$k \ge 3$ (Karp). Dado un grafo $G$, se genera el programa $P_G$
(\tm{lcero/reductions.py}), que tiene una variable \tm{v}$_i$ por vértice y
un \tm{if} por cada arista:

\begin{Verbatim}
int sel = 0;
int v0; int v1; ...
if (sel == 0) {        // arista {u, w}
    vu = 1;
    vw = 2;
    print(vu + vw);
} else if (sel == 1) { ...
\end{Verbatim}

En la rama de la arista $\{u, w\}$, cuando se asigna \tm{vw} la variable
\tm{vu} está viva (se usa en la suma), así que interfieren. Fuera de esa
rama ninguna de las dos está viva, porque cada rama las vuelve a asignar
antes de usarlas. Por lo tanto, entre las variables \tm{v}$_i$ el grafo de
interferencia es exactamente $G$. Las demás variables (\tm{sel} y los
temporales) no tocan a las \tm{v}$_i$ y forman un grafo que se colorea con
2 colores, así que para $k \ge 3$, $G$ se puede colorear con $k$ colores si
y solo si $P_G$ entra en $k$ registros. Como $P_G$ tiene un tamaño lineal en
el de $G$, la reducción es polinomial.
\end{proof}

Las pruebas verifican que el grafo de interferencia de $P_G$ es exactamente
$G$ para varios grafos conocidos (ciclos, grafos completos y el grafo de
Petersen) y para 30 grafos al azar.

\paragraph{Heurística de Chaitin}

Como el problema es NP-completo, el módulo \tm{lcero/regalloc.py} usa la
heurística de Chaitin:

\begin{enumerate}
  \item Calcular la vivacidad y armar el grafo de interferencia.
  \item \textbf{Simplificar}: mientras haya un nodo con menos de $k$
  vecinos, sacarlo del grafo y apilarlo. Ese nodo siempre se va a poder
  colorear, porque sus vecinos usan a lo sumo $k-1$ colores.
  \item Si todos los nodos restantes tienen $k$ o más vecinos,
  \textbf{elegir uno para mandar a memoria} (el de menor costo/grado, donde
  el costo cuenta sus usos y pesa 10 veces más por cada nivel de
  \tm{while}) y sacarlo.
  \item \textbf{Seleccionar}: desapilar y darle a cada nodo el primer color
  que no usen sus vecinos.
  \item Si se mandó algo a memoria, agregar un \tm{load} antes de cada uso y
  un \tm{store} después de cada definición, y volver al paso 1.
\end{enumerate}

\paragraph{Por qué una heurística y no una búsqueda exhaustiva}

Buscar un coloreo exacto puede requerir probar hasta $k^{V}$ asignaciones.
Se comparó Chaitin con una implementación exacta por backtracking sobre
grafos al azar (figura~\ref{fig:coloreo}): el tiempo del método exacto crece
como $\ColoreoBase^{V}$ (se duplica cada \ColoreoDuplica{} vértices),
mientras que Chaitin tarda décimas de milisegundo. La desventaja es que a
veces manda a memoria variables sin necesidad: el grafo de Petersen se puede
colorear con 3 colores, pero todos sus vértices tienen grado 3, así que con
$k = 3$ Chaitin no puede simplificar ninguno y manda uno a memoria.

\begin{figure}[H]
\centering
\includegraphics[width=0.85\textwidth]{generado/coloreo_exacto.pdf}
\caption{Coloreo exacto (backtracking) contra la heurística de Chaitin en grafos al azar.}
\label{fig:coloreo}
\end{figure}

La tabla~\ref{tab:registros} muestra los resultados sobre los programas de
prueba. Por ejemplo, en \tm{fibonacci} hay cuatro variables vivas a la vez
(\tm{a}, \tm{b}, \tm{i} y \tm{t}), así que con 3 registros no entran todas
y Chaitin manda dos a memoria (\tm{\FibSlots}).

\begin{table}[H]
\centering
\begin{grilla}
\begin{tabular}{|l|c|c|c|c|c|c|c|}
\hline
\multicolumn{5}{|c|}{} & \multicolumn{3}{c|}{\textbf{Variables a memoria (rondas)}} \\ \cline{6-8}
\thead{Programa} & \thead{$N$} & \thead{$V$} & \thead{$E$} & \thead{Grado\\máximo} & \thead{$k=2$} & \thead{$k=3$} & \thead{$k=4$} \\ \hline
\tabla{generado/f8_resultados_grid.tex}
\end{tabular}
\end{grilla}
\caption{Asignación de registros en los programas de prueba ($N$, $V$, $E$ y grado de la primera ronda).}
\label{tab:registros}
\end{table}

% =============================================================================
\section{Pruebas automatizadas}

El proyecto tiene \NTests{} pruebas automáticas (\tm{python -m pytest}). La
base son \NProgramas{} programas L-Cero en \tm{tests/programas/}, cada uno
con lo que se espera de él indicado en la primera línea
(\tm{// SALIDA: \dots}, \tm{// ERROR: fase línea:columna} o \tm{// BUCLE}).
La tabla~\ref{tab:pruebas} resume las categorías.

\begin{table}[H]
\centering
\begin{grilla}
\begin{tabular}{|L{2.9cm}|C{1.9cm}|L{8.9cm}|}
\hline
\thead{Categoría} & \thead{Programas} & \thead{Qué se verifica} \\ \hline
Válidos & \NProgValidos & La salida exacta, con 2, 3, 4 y 8 registros (factorial, Fibonacci, MCD, primos, Collatz, una MT, etc.). \\ \hline
Error léxico & \NProgErrorlexico & Que falle el lexer en la línea y columna indicadas (\tm{@}, \tm{ñ}, comillas, \tm{\&\&}, \tm{!}). \\ \hline
Error sintáctico & \NProgErrorsintactico & Que falle el parser donde corresponde (falta \tm{;}, llave sin cerrar, \tm{else} sin \tm{if}, etc.). \\ \hline
Error semántico & \NProgErrorsemantico & Doble declaración, variable sin declarar, sin inicializar o fuera de ámbito. \\ \hline
Bucles infinitos & \NProgBucleinfinito & Que compilen y no terminen en 200\,000 pasos. \\ \hline
\end{tabular}
\end{grilla}
\caption{Categorías de programas de prueba.}
\label{tab:pruebas}
\end{table}

Además hay pruebas de cada parte por separado: los First y Follow contra los
valores calculados a mano, la minimalidad del AFD, cada regla semántica, la
recuperación de errores del parser y del análisis semántico, los avisos del
compilador, la función de transición de la máquina, las dos reducciones y la
validez de la asignación de registros. El compilador se usa desde la línea
de comandos:

\begin{Verbatim}
python -m lcero ejecutar programa.l0 -k 4     # compila y ejecuta con 4 registros
python -m lcero avisos programa.l0            # solo los avisos del compilador
python -m lcero tokens programa.l0            # también: ast, ri, maquina, traza
python -m lcero registros programa.l0 -k 3    # grafo de interferencia y asignación
python -m lcero afd                           # matriz del AFD (también: gramatica)
\end{Verbatim}

\textbf{Repositorio:} \Repositorio

% =============================================================================
\section{Conclusiones}

El trabajo permitió ver cómo cada parte del compilador corresponde a un tema
de la materia. Los tokens son lenguajes regulares y alcanza con un AFD; la
sintaxis necesita una gramática libre de contexto y un autómata de pila; los
controles sobre las variables ya no son libres de contexto y requieren algo
equivalente a un autómata linealmente acotado; y la ejecución es
Turing-completa, lo que hace imposible decidir en general si un programa
termina.

En la parte de complejidad, la asignación de registros mostró por qué en la
práctica se usan heurísticas: el problema exacto es NP-completo y la
búsqueda exhaustiva crece exponencialmente, mientras que Chaitin es casi
lineal, aunque no siempre encuentre la mejor solución. Por último, el
profiling permitió detectar un costo cuadrático que en el análisis teórico
inicial había pasado inadvertido.

% =============================================================================
\section{Referencias}

\begin{enumerate}[label={[\arabic*]}, leftmargin=2em, itemsep=3pt]
\item A. Aho, M. Lam, R. Sethi y J. Ullman. \emph{Compiladores: principios, técnicas y herramientas}, 2.ª ed. Pearson, 2008.
\item J. Hopcroft, R. Motwani y J. Ullman. \emph{Introducción a la teoría de autómatas, lenguajes y computación}, 3.ª ed. Pearson, 2007.
\item M. Sipser. \emph{Introduction to the Theory of Computation}, 3.ª ed. Cengage, 2013.
\item G. Chaitin et al. Register allocation via coloring. \emph{Computer Languages}, 6(1):47--57, 1981.
\item J. Hopcroft. An $n \log n$ algorithm for minimizing states in a finite automaton. En \emph{Theory of Machines and Computations}, Academic Press, 1971.
\item R. Karp. Reducibility among combinatorial problems. En \emph{Complexity of Computer Computations}, Plenum, 1972.
\item M. Minsky. \emph{Computation: Finite and Infinite Machines}. Prentice-Hall, 1967.
\item K. Thompson. Regular expression search algorithm. \emph{Communications of the ACM}, 11(6):419--422, 1968.
\item A. Turing. On computable numbers, with an application to the Entscheidungsproblem. \emph{Proc. London Math. Society}, 42:230--265, 1936.
\end{enumerate}

% =============================================================================
\newpage
\section*{Anexo A. AFD mínimo completo}
\addcontentsline{toc}{section}{Anexo A. AFD mínimo completo}

\begin{figure}[H]
\centering
\includegraphics[height=0.74\textheight]{generado/afd_simbolos.pdf}
\caption{Parte del AFD mínimo que reconoce operadores, delimitadores, comentarios y espacios.}
\end{figure}

\newpage
La matriz de transición $\delta$ tiene \NMin{} estados y \NClases{}
columnas. $\rightarrow$ indica el estado inicial, $\ast$ los estados finales
y $\cdot$ una transición al sumidero $q_\emptyset$. Columnas:
\tm{e f h i l n p r s t w}, una por cada letra que aparece en una palabra
reservada; \tm{L}, el resto de las letras y el \tm{\_}; \tm{D}, los dígitos;
un símbolo por columna; \tm{esp}, espacio, tabulación y
\tm{\textbackslash r}; \tm{\textbackslash n}; y \tm{otro}, cualquier otro
carácter (solo válido dentro de un comentario).

\begingroup
\scriptsize
\setlength{\tabcolsep}{1.6pt}
\renewcommand{\arraystretch}{1.12}
\begin{longtable}{@{}l|*{30}{c}@{}}
\toprule
$\delta$ & \texttt{e} & \texttt{f} & \texttt{h} & \texttt{i} & \texttt{l} & \texttt{n} & \texttt{p} & \texttt{r} & \texttt{s} & \texttt{t} & \texttt{w} & \texttt{L} & \texttt{D} & \texttt{+} & \texttt{-} & \texttt{*} & \texttt{/} & \texttt{\%} & \texttt{=} & \texttt{!} & \texttt{<} & \texttt{>} & \texttt{(} & \texttt{)} & \texttt{\{} & \texttt{\}} & \texttt{;} & \texttt{esp} & \texttt{\textbackslash{}n} & \texttt{otro} \\
\midrule
\endhead
\bottomrule
\endfoot
$\rightarrow$$q_{0}$ & 1 & 2 & 2 & 3 & 2 & 2 & 4 & 2 & 2 & 2 & 5 & 2 & 6 & 7 & 8 & 9 & 10 & 11 & 12 & 13 & 14 & 15 & 16 & 17 & 18 & 19 & 20 & 21 & 21 & $\cdot$ \\
$\ast$$q_{1}$ & 2 & 2 & 2 & 2 & 22 & 2 & 2 & 2 & 2 & 2 & 2 & 2 & 2 & $\cdot$ & $\cdot$ & $\cdot$ & $\cdot$ & $\cdot$ & $\cdot$ & $\cdot$ & $\cdot$ & $\cdot$ & $\cdot$ & $\cdot$ & $\cdot$ & $\cdot$ & $\cdot$ & $\cdot$ & $\cdot$ & $\cdot$ \\
$\ast$$q_{2}$ & 2 & 2 & 2 & 2 & 2 & 2 & 2 & 2 & 2 & 2 & 2 & 2 & 2 & $\cdot$ & $\cdot$ & $\cdot$ & $\cdot$ & $\cdot$ & $\cdot$ & $\cdot$ & $\cdot$ & $\cdot$ & $\cdot$ & $\cdot$ & $\cdot$ & $\cdot$ & $\cdot$ & $\cdot$ & $\cdot$ & $\cdot$ \\
$\ast$$q_{3}$ & 2 & 23 & 2 & 2 & 2 & 24 & 2 & 2 & 2 & 2 & 2 & 2 & 2 & $\cdot$ & $\cdot$ & $\cdot$ & $\cdot$ & $\cdot$ & $\cdot$ & $\cdot$ & $\cdot$ & $\cdot$ & $\cdot$ & $\cdot$ & $\cdot$ & $\cdot$ & $\cdot$ & $\cdot$ & $\cdot$ & $\cdot$ \\
$\ast$$q_{4}$ & 2 & 2 & 2 & 2 & 2 & 2 & 2 & 25 & 2 & 2 & 2 & 2 & 2 & $\cdot$ & $\cdot$ & $\cdot$ & $\cdot$ & $\cdot$ & $\cdot$ & $\cdot$ & $\cdot$ & $\cdot$ & $\cdot$ & $\cdot$ & $\cdot$ & $\cdot$ & $\cdot$ & $\cdot$ & $\cdot$ & $\cdot$ \\
$\ast$$q_{5}$ & 2 & 2 & 26 & 2 & 2 & 2 & 2 & 2 & 2 & 2 & 2 & 2 & 2 & $\cdot$ & $\cdot$ & $\cdot$ & $\cdot$ & $\cdot$ & $\cdot$ & $\cdot$ & $\cdot$ & $\cdot$ & $\cdot$ & $\cdot$ & $\cdot$ & $\cdot$ & $\cdot$ & $\cdot$ & $\cdot$ & $\cdot$ \\
$\ast$$q_{6}$ & $\cdot$ & $\cdot$ & $\cdot$ & $\cdot$ & $\cdot$ & $\cdot$ & $\cdot$ & $\cdot$ & $\cdot$ & $\cdot$ & $\cdot$ & $\cdot$ & 6 & $\cdot$ & $\cdot$ & $\cdot$ & $\cdot$ & $\cdot$ & $\cdot$ & $\cdot$ & $\cdot$ & $\cdot$ & $\cdot$ & $\cdot$ & $\cdot$ & $\cdot$ & $\cdot$ & $\cdot$ & $\cdot$ & $\cdot$ \\
$\ast$$q_{7}$ & $\cdot$ & $\cdot$ & $\cdot$ & $\cdot$ & $\cdot$ & $\cdot$ & $\cdot$ & $\cdot$ & $\cdot$ & $\cdot$ & $\cdot$ & $\cdot$ & $\cdot$ & $\cdot$ & $\cdot$ & $\cdot$ & $\cdot$ & $\cdot$ & $\cdot$ & $\cdot$ & $\cdot$ & $\cdot$ & $\cdot$ & $\cdot$ & $\cdot$ & $\cdot$ & $\cdot$ & $\cdot$ & $\cdot$ & $\cdot$ \\
$\ast$$q_{8}$ & $\cdot$ & $\cdot$ & $\cdot$ & $\cdot$ & $\cdot$ & $\cdot$ & $\cdot$ & $\cdot$ & $\cdot$ & $\cdot$ & $\cdot$ & $\cdot$ & $\cdot$ & $\cdot$ & $\cdot$ & $\cdot$ & $\cdot$ & $\cdot$ & $\cdot$ & $\cdot$ & $\cdot$ & $\cdot$ & $\cdot$ & $\cdot$ & $\cdot$ & $\cdot$ & $\cdot$ & $\cdot$ & $\cdot$ & $\cdot$ \\
$\ast$$q_{9}$ & $\cdot$ & $\cdot$ & $\cdot$ & $\cdot$ & $\cdot$ & $\cdot$ & $\cdot$ & $\cdot$ & $\cdot$ & $\cdot$ & $\cdot$ & $\cdot$ & $\cdot$ & $\cdot$ & $\cdot$ & $\cdot$ & $\cdot$ & $\cdot$ & $\cdot$ & $\cdot$ & $\cdot$ & $\cdot$ & $\cdot$ & $\cdot$ & $\cdot$ & $\cdot$ & $\cdot$ & $\cdot$ & $\cdot$ & $\cdot$ \\
$\ast$$q_{10}$ & $\cdot$ & $\cdot$ & $\cdot$ & $\cdot$ & $\cdot$ & $\cdot$ & $\cdot$ & $\cdot$ & $\cdot$ & $\cdot$ & $\cdot$ & $\cdot$ & $\cdot$ & $\cdot$ & $\cdot$ & $\cdot$ & 27 & $\cdot$ & $\cdot$ & $\cdot$ & $\cdot$ & $\cdot$ & $\cdot$ & $\cdot$ & $\cdot$ & $\cdot$ & $\cdot$ & $\cdot$ & $\cdot$ & $\cdot$ \\
$\ast$$q_{11}$ & $\cdot$ & $\cdot$ & $\cdot$ & $\cdot$ & $\cdot$ & $\cdot$ & $\cdot$ & $\cdot$ & $\cdot$ & $\cdot$ & $\cdot$ & $\cdot$ & $\cdot$ & $\cdot$ & $\cdot$ & $\cdot$ & $\cdot$ & $\cdot$ & $\cdot$ & $\cdot$ & $\cdot$ & $\cdot$ & $\cdot$ & $\cdot$ & $\cdot$ & $\cdot$ & $\cdot$ & $\cdot$ & $\cdot$ & $\cdot$ \\
$\ast$$q_{12}$ & $\cdot$ & $\cdot$ & $\cdot$ & $\cdot$ & $\cdot$ & $\cdot$ & $\cdot$ & $\cdot$ & $\cdot$ & $\cdot$ & $\cdot$ & $\cdot$ & $\cdot$ & $\cdot$ & $\cdot$ & $\cdot$ & $\cdot$ & $\cdot$ & 28 & $\cdot$ & $\cdot$ & $\cdot$ & $\cdot$ & $\cdot$ & $\cdot$ & $\cdot$ & $\cdot$ & $\cdot$ & $\cdot$ & $\cdot$ \\
$q_{13}$ & $\cdot$ & $\cdot$ & $\cdot$ & $\cdot$ & $\cdot$ & $\cdot$ & $\cdot$ & $\cdot$ & $\cdot$ & $\cdot$ & $\cdot$ & $\cdot$ & $\cdot$ & $\cdot$ & $\cdot$ & $\cdot$ & $\cdot$ & $\cdot$ & 29 & $\cdot$ & $\cdot$ & $\cdot$ & $\cdot$ & $\cdot$ & $\cdot$ & $\cdot$ & $\cdot$ & $\cdot$ & $\cdot$ & $\cdot$ \\
$\ast$$q_{14}$ & $\cdot$ & $\cdot$ & $\cdot$ & $\cdot$ & $\cdot$ & $\cdot$ & $\cdot$ & $\cdot$ & $\cdot$ & $\cdot$ & $\cdot$ & $\cdot$ & $\cdot$ & $\cdot$ & $\cdot$ & $\cdot$ & $\cdot$ & $\cdot$ & 30 & $\cdot$ & $\cdot$ & $\cdot$ & $\cdot$ & $\cdot$ & $\cdot$ & $\cdot$ & $\cdot$ & $\cdot$ & $\cdot$ & $\cdot$ \\
$\ast$$q_{15}$ & $\cdot$ & $\cdot$ & $\cdot$ & $\cdot$ & $\cdot$ & $\cdot$ & $\cdot$ & $\cdot$ & $\cdot$ & $\cdot$ & $\cdot$ & $\cdot$ & $\cdot$ & $\cdot$ & $\cdot$ & $\cdot$ & $\cdot$ & $\cdot$ & 31 & $\cdot$ & $\cdot$ & $\cdot$ & $\cdot$ & $\cdot$ & $\cdot$ & $\cdot$ & $\cdot$ & $\cdot$ & $\cdot$ & $\cdot$ \\
$\ast$$q_{16}$ & $\cdot$ & $\cdot$ & $\cdot$ & $\cdot$ & $\cdot$ & $\cdot$ & $\cdot$ & $\cdot$ & $\cdot$ & $\cdot$ & $\cdot$ & $\cdot$ & $\cdot$ & $\cdot$ & $\cdot$ & $\cdot$ & $\cdot$ & $\cdot$ & $\cdot$ & $\cdot$ & $\cdot$ & $\cdot$ & $\cdot$ & $\cdot$ & $\cdot$ & $\cdot$ & $\cdot$ & $\cdot$ & $\cdot$ & $\cdot$ \\
$\ast$$q_{17}$ & $\cdot$ & $\cdot$ & $\cdot$ & $\cdot$ & $\cdot$ & $\cdot$ & $\cdot$ & $\cdot$ & $\cdot$ & $\cdot$ & $\cdot$ & $\cdot$ & $\cdot$ & $\cdot$ & $\cdot$ & $\cdot$ & $\cdot$ & $\cdot$ & $\cdot$ & $\cdot$ & $\cdot$ & $\cdot$ & $\cdot$ & $\cdot$ & $\cdot$ & $\cdot$ & $\cdot$ & $\cdot$ & $\cdot$ & $\cdot$ \\
$\ast$$q_{18}$ & $\cdot$ & $\cdot$ & $\cdot$ & $\cdot$ & $\cdot$ & $\cdot$ & $\cdot$ & $\cdot$ & $\cdot$ & $\cdot$ & $\cdot$ & $\cdot$ & $\cdot$ & $\cdot$ & $\cdot$ & $\cdot$ & $\cdot$ & $\cdot$ & $\cdot$ & $\cdot$ & $\cdot$ & $\cdot$ & $\cdot$ & $\cdot$ & $\cdot$ & $\cdot$ & $\cdot$ & $\cdot$ & $\cdot$ & $\cdot$ \\
$\ast$$q_{19}$ & $\cdot$ & $\cdot$ & $\cdot$ & $\cdot$ & $\cdot$ & $\cdot$ & $\cdot$ & $\cdot$ & $\cdot$ & $\cdot$ & $\cdot$ & $\cdot$ & $\cdot$ & $\cdot$ & $\cdot$ & $\cdot$ & $\cdot$ & $\cdot$ & $\cdot$ & $\cdot$ & $\cdot$ & $\cdot$ & $\cdot$ & $\cdot$ & $\cdot$ & $\cdot$ & $\cdot$ & $\cdot$ & $\cdot$ & $\cdot$ \\
$\ast$$q_{20}$ & $\cdot$ & $\cdot$ & $\cdot$ & $\cdot$ & $\cdot$ & $\cdot$ & $\cdot$ & $\cdot$ & $\cdot$ & $\cdot$ & $\cdot$ & $\cdot$ & $\cdot$ & $\cdot$ & $\cdot$ & $\cdot$ & $\cdot$ & $\cdot$ & $\cdot$ & $\cdot$ & $\cdot$ & $\cdot$ & $\cdot$ & $\cdot$ & $\cdot$ & $\cdot$ & $\cdot$ & $\cdot$ & $\cdot$ & $\cdot$ \\
$\ast$$q_{21}$ & $\cdot$ & $\cdot$ & $\cdot$ & $\cdot$ & $\cdot$ & $\cdot$ & $\cdot$ & $\cdot$ & $\cdot$ & $\cdot$ & $\cdot$ & $\cdot$ & $\cdot$ & $\cdot$ & $\cdot$ & $\cdot$ & $\cdot$ & $\cdot$ & $\cdot$ & $\cdot$ & $\cdot$ & $\cdot$ & $\cdot$ & $\cdot$ & $\cdot$ & $\cdot$ & $\cdot$ & 21 & 21 & $\cdot$ \\
$\ast$$q_{22}$ & 2 & 2 & 2 & 2 & 2 & 2 & 2 & 2 & 32 & 2 & 2 & 2 & 2 & $\cdot$ & $\cdot$ & $\cdot$ & $\cdot$ & $\cdot$ & $\cdot$ & $\cdot$ & $\cdot$ & $\cdot$ & $\cdot$ & $\cdot$ & $\cdot$ & $\cdot$ & $\cdot$ & $\cdot$ & $\cdot$ & $\cdot$ \\
$\ast$$q_{23}$ & 2 & 2 & 2 & 2 & 2 & 2 & 2 & 2 & 2 & 2 & 2 & 2 & 2 & $\cdot$ & $\cdot$ & $\cdot$ & $\cdot$ & $\cdot$ & $\cdot$ & $\cdot$ & $\cdot$ & $\cdot$ & $\cdot$ & $\cdot$ & $\cdot$ & $\cdot$ & $\cdot$ & $\cdot$ & $\cdot$ & $\cdot$ \\
$\ast$$q_{24}$ & 2 & 2 & 2 & 2 & 2 & 2 & 2 & 2 & 2 & 33 & 2 & 2 & 2 & $\cdot$ & $\cdot$ & $\cdot$ & $\cdot$ & $\cdot$ & $\cdot$ & $\cdot$ & $\cdot$ & $\cdot$ & $\cdot$ & $\cdot$ & $\cdot$ & $\cdot$ & $\cdot$ & $\cdot$ & $\cdot$ & $\cdot$ \\
$\ast$$q_{25}$ & 2 & 2 & 2 & 34 & 2 & 2 & 2 & 2 & 2 & 2 & 2 & 2 & 2 & $\cdot$ & $\cdot$ & $\cdot$ & $\cdot$ & $\cdot$ & $\cdot$ & $\cdot$ & $\cdot$ & $\cdot$ & $\cdot$ & $\cdot$ & $\cdot$ & $\cdot$ & $\cdot$ & $\cdot$ & $\cdot$ & $\cdot$ \\
$\ast$$q_{26}$ & 2 & 2 & 2 & 35 & 2 & 2 & 2 & 2 & 2 & 2 & 2 & 2 & 2 & $\cdot$ & $\cdot$ & $\cdot$ & $\cdot$ & $\cdot$ & $\cdot$ & $\cdot$ & $\cdot$ & $\cdot$ & $\cdot$ & $\cdot$ & $\cdot$ & $\cdot$ & $\cdot$ & $\cdot$ & $\cdot$ & $\cdot$ \\
$\ast$$q_{27}$ & 27 & 27 & 27 & 27 & 27 & 27 & 27 & 27 & 27 & 27 & 27 & 27 & 27 & 27 & 27 & 27 & 27 & 27 & 27 & 27 & 27 & 27 & 27 & 27 & 27 & 27 & 27 & 27 & $\cdot$ & 27 \\
$\ast$$q_{28}$ & $\cdot$ & $\cdot$ & $\cdot$ & $\cdot$ & $\cdot$ & $\cdot$ & $\cdot$ & $\cdot$ & $\cdot$ & $\cdot$ & $\cdot$ & $\cdot$ & $\cdot$ & $\cdot$ & $\cdot$ & $\cdot$ & $\cdot$ & $\cdot$ & $\cdot$ & $\cdot$ & $\cdot$ & $\cdot$ & $\cdot$ & $\cdot$ & $\cdot$ & $\cdot$ & $\cdot$ & $\cdot$ & $\cdot$ & $\cdot$ \\
$\ast$$q_{29}$ & $\cdot$ & $\cdot$ & $\cdot$ & $\cdot$ & $\cdot$ & $\cdot$ & $\cdot$ & $\cdot$ & $\cdot$ & $\cdot$ & $\cdot$ & $\cdot$ & $\cdot$ & $\cdot$ & $\cdot$ & $\cdot$ & $\cdot$ & $\cdot$ & $\cdot$ & $\cdot$ & $\cdot$ & $\cdot$ & $\cdot$ & $\cdot$ & $\cdot$ & $\cdot$ & $\cdot$ & $\cdot$ & $\cdot$ & $\cdot$ \\
$\ast$$q_{30}$ & $\cdot$ & $\cdot$ & $\cdot$ & $\cdot$ & $\cdot$ & $\cdot$ & $\cdot$ & $\cdot$ & $\cdot$ & $\cdot$ & $\cdot$ & $\cdot$ & $\cdot$ & $\cdot$ & $\cdot$ & $\cdot$ & $\cdot$ & $\cdot$ & $\cdot$ & $\cdot$ & $\cdot$ & $\cdot$ & $\cdot$ & $\cdot$ & $\cdot$ & $\cdot$ & $\cdot$ & $\cdot$ & $\cdot$ & $\cdot$ \\
$\ast$$q_{31}$ & $\cdot$ & $\cdot$ & $\cdot$ & $\cdot$ & $\cdot$ & $\cdot$ & $\cdot$ & $\cdot$ & $\cdot$ & $\cdot$ & $\cdot$ & $\cdot$ & $\cdot$ & $\cdot$ & $\cdot$ & $\cdot$ & $\cdot$ & $\cdot$ & $\cdot$ & $\cdot$ & $\cdot$ & $\cdot$ & $\cdot$ & $\cdot$ & $\cdot$ & $\cdot$ & $\cdot$ & $\cdot$ & $\cdot$ & $\cdot$ \\
$\ast$$q_{32}$ & 36 & 2 & 2 & 2 & 2 & 2 & 2 & 2 & 2 & 2 & 2 & 2 & 2 & $\cdot$ & $\cdot$ & $\cdot$ & $\cdot$ & $\cdot$ & $\cdot$ & $\cdot$ & $\cdot$ & $\cdot$ & $\cdot$ & $\cdot$ & $\cdot$ & $\cdot$ & $\cdot$ & $\cdot$ & $\cdot$ & $\cdot$ \\
$\ast$$q_{33}$ & 2 & 2 & 2 & 2 & 2 & 2 & 2 & 2 & 2 & 2 & 2 & 2 & 2 & $\cdot$ & $\cdot$ & $\cdot$ & $\cdot$ & $\cdot$ & $\cdot$ & $\cdot$ & $\cdot$ & $\cdot$ & $\cdot$ & $\cdot$ & $\cdot$ & $\cdot$ & $\cdot$ & $\cdot$ & $\cdot$ & $\cdot$ \\
$\ast$$q_{34}$ & 2 & 2 & 2 & 2 & 2 & 37 & 2 & 2 & 2 & 2 & 2 & 2 & 2 & $\cdot$ & $\cdot$ & $\cdot$ & $\cdot$ & $\cdot$ & $\cdot$ & $\cdot$ & $\cdot$ & $\cdot$ & $\cdot$ & $\cdot$ & $\cdot$ & $\cdot$ & $\cdot$ & $\cdot$ & $\cdot$ & $\cdot$ \\
$\ast$$q_{35}$ & 2 & 2 & 2 & 2 & 38 & 2 & 2 & 2 & 2 & 2 & 2 & 2 & 2 & $\cdot$ & $\cdot$ & $\cdot$ & $\cdot$ & $\cdot$ & $\cdot$ & $\cdot$ & $\cdot$ & $\cdot$ & $\cdot$ & $\cdot$ & $\cdot$ & $\cdot$ & $\cdot$ & $\cdot$ & $\cdot$ & $\cdot$ \\
$\ast$$q_{36}$ & 2 & 2 & 2 & 2 & 2 & 2 & 2 & 2 & 2 & 2 & 2 & 2 & 2 & $\cdot$ & $\cdot$ & $\cdot$ & $\cdot$ & $\cdot$ & $\cdot$ & $\cdot$ & $\cdot$ & $\cdot$ & $\cdot$ & $\cdot$ & $\cdot$ & $\cdot$ & $\cdot$ & $\cdot$ & $\cdot$ & $\cdot$ \\
$\ast$$q_{37}$ & 2 & 2 & 2 & 2 & 2 & 2 & 2 & 2 & 2 & 39 & 2 & 2 & 2 & $\cdot$ & $\cdot$ & $\cdot$ & $\cdot$ & $\cdot$ & $\cdot$ & $\cdot$ & $\cdot$ & $\cdot$ & $\cdot$ & $\cdot$ & $\cdot$ & $\cdot$ & $\cdot$ & $\cdot$ & $\cdot$ & $\cdot$ \\
$\ast$$q_{38}$ & 40 & 2 & 2 & 2 & 2 & 2 & 2 & 2 & 2 & 2 & 2 & 2 & 2 & $\cdot$ & $\cdot$ & $\cdot$ & $\cdot$ & $\cdot$ & $\cdot$ & $\cdot$ & $\cdot$ & $\cdot$ & $\cdot$ & $\cdot$ & $\cdot$ & $\cdot$ & $\cdot$ & $\cdot$ & $\cdot$ & $\cdot$ \\
$\ast$$q_{39}$ & 2 & 2 & 2 & 2 & 2 & 2 & 2 & 2 & 2 & 2 & 2 & 2 & 2 & $\cdot$ & $\cdot$ & $\cdot$ & $\cdot$ & $\cdot$ & $\cdot$ & $\cdot$ & $\cdot$ & $\cdot$ & $\cdot$ & $\cdot$ & $\cdot$ & $\cdot$ & $\cdot$ & $\cdot$ & $\cdot$ & $\cdot$ \\
$\ast$$q_{40}$ & 2 & 2 & 2 & 2 & 2 & 2 & 2 & 2 & 2 & 2 & 2 & 2 & 2 & $\cdot$ & $\cdot$ & $\cdot$ & $\cdot$ & $\cdot$ & $\cdot$ & $\cdot$ & $\cdot$ & $\cdot$ & $\cdot$ & $\cdot$ & $\cdot$ & $\cdot$ & $\cdot$ & $\cdot$ & $\cdot$ & $\cdot$ \\
$q_\emptyset$ & $\cdot$ & $\cdot$ & $\cdot$ & $\cdot$ & $\cdot$ & $\cdot$ & $\cdot$ & $\cdot$ & $\cdot$ & $\cdot$ & $\cdot$ & $\cdot$ & $\cdot$ & $\cdot$ & $\cdot$ & $\cdot$ & $\cdot$ & $\cdot$ & $\cdot$ & $\cdot$ & $\cdot$ & $\cdot$ & $\cdot$ & $\cdot$ & $\cdot$ & $\cdot$ & $\cdot$ & $\cdot$ & $\cdot$ & $\cdot$ \\
\end{longtable}

\endgroup

\newpage
\section*{Anexo B. Cálculo de First y Follow}
\addcontentsline{toc}{section}{Anexo B. Cálculo de First y Follow}

\begingroup\footnotesize\renewcommand{\arraystretch}{1.2}
\begin{longtable}{|l|l|l|}
\caption{First: aporte de cada producción.}\\
\hline
\thead{$A$} & \thead{Producción} & \thead{$\First$ del lado derecho} \\ \hline
\endfirsthead
\hline
\thead{$A$} & \thead{Producción} & \thead{$\First$ del lado derecho} \\ \hline
\endhead
\tabla{generado/f2_first_calculo.tex}
\hline
\end{longtable}
\endgroup

\begingroup\footnotesize\renewcommand{\arraystretch}{1.2}
\begin{longtable}{|l|l|p{6.3cm}|}
\caption{Follow: reglas aplicadas (entre paréntesis, la producción donde aparece cada caso).}\\
\hline
\thead{$B$} & \thead{Regla} & \thead{Aporte} \\ \hline
\endfirsthead
\hline
\thead{$B$} & \thead{Regla} & \thead{Aporte} \\ \hline
\endhead
\tabla{generado/f2_follow_calculo.tex}
\hline
\end{longtable}
\endgroup

\section*{Anexo C. Tabla LL(1)}
\addcontentsline{toc}{section}{Anexo C. Tabla LL(1)}

\begin{table}[H]
\centering
\scriptsize
\setlength{\tabcolsep}{2.3pt}
\begin{tabular}{@{}l|*{25}{c}@{}}
\toprule
 & \rotatebox{90}{\tm{int}} & \rotatebox{90}{\tm{if}} & \rotatebox{90}{\tm{else}} & \rotatebox{90}{\tm{while}} & \rotatebox{90}{\tm{print}} & \rotatebox{90}{\tm{id}} & \rotatebox{90}{\tm{num}} & \rotatebox{90}{\tm{+}} & \rotatebox{90}{\tm{-}} & \rotatebox{90}{\tm{*}} & \rotatebox{90}{\tm{/}} & \rotatebox{90}{\tm{\%}} & \rotatebox{90}{\tm{=}} & \rotatebox{90}{\tm{==}} & \rotatebox{90}{\tm{!=}} & \rotatebox{90}{\tm{<}} & \rotatebox{90}{\tm{<=}} & \rotatebox{90}{\tm{>}} & \rotatebox{90}{\tm{>=}} & \rotatebox{90}{\tm{(}} & \rotatebox{90}{\tm{)}} & \rotatebox{90}{\tm{\{}} & \rotatebox{90}{\tm{\}}} & \rotatebox{90}{\tm{;}} & \rotatebox{90}{\tm{\$}} \\
\midrule
\nt{Programa} & 1 & 1 &  & 1 & 1 & 1 &  &  &  &  &  &  &  &  &  &  &  &  &  &  &  & 1 &  &  & 1 \\
\nt{ListaSent} & 2 & 2 &  & 2 & 2 & 2 &  &  &  &  &  &  &  &  &  &  &  &  &  &  &  & 2 & 3 &  & 3 \\
\nt{Sent} & 4 & 6 &  & 7 & 8 & 5 &  &  &  &  &  &  &  &  &  &  &  &  &  &  &  & 9 &  &  &  \\
\nt{Decl} & 10 &  &  &  &  &  &  &  &  &  &  &  &  &  &  &  &  &  &  &  &  &  &  &  &  \\
\nt{DeclInit} &  &  &  &  &  &  &  &  &  &  &  &  & 11 &  &  &  &  &  &  &  &  &  &  & 12 &  \\
\nt{Asig} &  &  &  &  &  & 13 &  &  &  &  &  &  &  &  &  &  &  &  &  &  &  &  &  &  &  \\
\nt{Si} &  & 14 &  &  &  &  &  &  &  &  &  &  &  &  &  &  &  &  &  &  &  &  &  &  &  \\
\nt{SiNo} & 16 & 16 & 15 & 16 & 16 & 16 &  &  &  &  &  &  &  &  &  &  &  &  &  &  &  & 16 & 16 &  & 16 \\
\nt{CuerpoSiNo} &  & 18 &  &  &  &  &  &  &  &  &  &  &  &  &  &  &  &  &  &  &  & 17 &  &  &  \\
\nt{Mientras} &  &  &  & 19 &  &  &  &  &  &  &  &  &  &  &  &  &  &  &  &  &  &  &  &  &  \\
\nt{Imprimir} &  &  &  &  & 20 &  &  &  &  &  &  &  &  &  &  &  &  &  &  &  &  &  &  &  &  \\
\nt{Bloque} &  &  &  &  &  &  &  &  &  &  &  &  &  &  &  &  &  &  &  &  &  & 21 &  &  &  \\
\nt{Expr} &  &  &  &  &  & 22 & 22 &  & 22 &  &  &  &  &  &  &  &  &  &  & 22 &  &  &  &  &  \\
\nt{ExprRel} &  &  &  &  &  &  &  &  &  &  &  &  &  & 23 & 23 & 23 & 23 & 23 & 23 &  & 24 &  &  & 24 &  \\
\nt{OpRel} &  &  &  &  &  &  &  &  &  &  &  &  &  & 25 & 26 & 27 & 28 & 29 & 30 &  &  &  &  &  &  \\
\nt{Arit} &  &  &  &  &  & 31 & 31 &  & 31 &  &  &  &  &  &  &  &  &  &  & 31 &  &  &  &  &  \\
\nt{AritCola} &  &  &  &  &  &  &  & 32 & 33 &  &  &  &  & 34 & 34 & 34 & 34 & 34 & 34 &  & 34 &  &  & 34 &  \\
\nt{Term} &  &  &  &  &  & 35 & 35 &  & 35 &  &  &  &  &  &  &  &  &  &  & 35 &  &  &  &  &  \\
\nt{TermCola} &  &  &  &  &  &  &  & 39 & 39 & 36 & 37 & 38 &  & 39 & 39 & 39 & 39 & 39 & 39 &  & 39 &  &  & 39 &  \\
\nt{Factor} &  &  &  &  &  & 41 & 40 &  & 43 &  &  &  &  &  &  &  &  &  &  & 42 &  &  &  &  &  \\
\bottomrule
\end{tabular}

\caption{Tabla $M[A, a]$: número de la producción que se aplica con $A$ en el tope y $a$ como próximo token. Las celdas vacías son errores.}
\end{table}

\end{document}
